\documentclass[11pt]{article}
\usepackage[margin=1in]{geometry}
\usepackage{amsmath,amssymb,amsthm}
\usepackage{graphicx}
\usepackage[hidelinks]{hyperref}

\newtheorem{theorem}{Theorem}
\newtheorem{corollary}{Corollary}

\newcommand{\ellTwo}{\ell_2}

\begin{document}

\begin{center}
{\Large Partial result for Godefroy's Problem 7}\\[0.4em]
{\large The non-\(\ell_2\)-saturated case has countably many non-isomorphic subspaces}\\[0.4em]
status: partial result, not a full solution
\end{center}

\section*{Source Problem}

The source is Dodos--Lopez-Abad--Todorcevic,
\emph{Banach spaces and Ramsey Theory: some open problems},
arXiv:1006.2668v1. Problem 7, attributed to Godefroy, asks:

\begin{quote}
Let \(X\) be a separable non-Hilbertian Banach space. Is it true that there
exists a sequence \((X_n)\) of subspaces of \(X\) which are pairwise
non-isomorphic?
\end{quote}

\begin{center}
\includegraphics[width=\textwidth]{figures/open_problem_crop.png}
\end{center}

\section*{Definitions}

A separable Banach space \(X\) is \(\ell_2\)-saturated if every closed
infinite-dimensional subspace of \(X\) contains a further subspace isomorphic
to \(\ell_2\).

We use the homogeneous Banach-space theorem of Komorowski--Tomczak-Jaegermann
and Gowers: if a separable infinite-dimensional Banach space is isomorphic to
all of its closed infinite-dimensional subspaces, then it is isomorphic to
\(\ell_2\).

\section*{Idea of the Proof}

The obstruction is finite rather than analytic. If \(X\) did not contain a
sequence of pairwise non-isomorphic subspaces, then the set of isomorphism
classes of subspaces of \(X\) would be finite. Inside any finite partially
ordered family of subspace-classes, one can choose a minimal class. A
representative \(Y\) of such a minimal class is isomorphic to each of its own
infinite-dimensional subspaces; otherwise one would move to a strictly smaller
class. The homogeneous Banach-space theorem then forces \(Y\simeq\ell_2\).
Repeating this inside every subspace shows that \(X\) must be \(\ell_2\)-saturated.

\section*{The Partial Result}

\begin{theorem}
Let \(X\) be a separable infinite-dimensional Banach space. If \(X\) is not
\(\ell_2\)-saturated, then \(X\) contains a sequence of closed
infinite-dimensional pairwise non-isomorphic subspaces unless \(X\) is already
outside the hypothesis by being Hilbertian. In particular, every separable
non-Hilbertian Banach space which is not \(\ell_2\)-saturated satisfies the
conclusion of Godefroy's Problem 7.
\end{theorem}

\begin{proof}
We prove the contrapositive. Suppose that \(X\) has no sequence of closed
infinite-dimensional pairwise non-isomorphic subspaces. Then the set
\(\mathcal C(X)\) of isomorphism classes of closed infinite-dimensional
subspaces of \(X\) is finite.

For classes \(A,B\in\mathcal C(X)\), write \(A\preceq B\) if a representative
of \(A\) is isomorphic to a closed infinite-dimensional subspace of a
representative of \(B\). This is a finite preorder. Choose a minimal class
\(A_0\) and let \(Y\subset X\) be a representative of \(A_0\). If \(Z\subset Y\)
is a closed infinite-dimensional subspace, then \([Z]\preceq [Y]=A_0\). By
minimality, \([Z]=A_0\). Thus \(Y\) is isomorphic to all of its closed
infinite-dimensional subspaces. By the homogeneous Banach-space theorem,
\(Y\) is isomorphic to \(\ell_2\).

Now let \(W\subset X\) be any closed infinite-dimensional subspace. The same
finite-class argument can be applied inside \(W\): among the isomorphism
classes represented by subspaces of \(W\), choose a minimal one. A representative
of that minimal class is isomorphic to all of its closed infinite-dimensional
subspaces, hence isomorphic to \(\ell_2\). Therefore every closed
infinite-dimensional subspace \(W\) of \(X\) contains a further subspace
isomorphic to \(\ell_2\). That is, \(X\) is \(\ell_2\)-saturated.

Consequently, if \(X\) is not \(\ell_2\)-saturated, the supposition was false:
\(X\) must contain a sequence of pairwise non-isomorphic subspaces.
\end{proof}

\begin{corollary}
Any counterexample to Godefroy's Problem 7 must be a separable non-Hilbertian
\(\ell_2\)-saturated Banach space.
\end{corollary}

\section*{Limitations and Novelty Check}

This is a reduction, not a full solution. It leaves open the \(\ell_2\)-saturated
case. It also gives only countably many pairwise non-isomorphic subspaces in the
non-\(\ell_2\)-saturated case; it does not prove the continuum-many strengthening
from Problem 8.

The run indexes were searched for arXiv:1006.2668, the survey title, Godefroy's
Problem 7, and pairwise non-isomorphic subspaces. A prior attempt for a different
part of the same survey, the Rosenthal/unique-spreading-model edge, was read and
not duplicated. Bounded web searches for the exact \(\ell_2\)-saturated reduction
did not reveal an obvious prior packet or later exact answer during this pass.

\section*{References}

\begin{thebibliography}{9}

\bibitem{DLT}
P. Dodos, J. Lopez-Abad, and S. Todorcevic,
\emph{Banach spaces and Ramsey Theory: some open problems},
arXiv:1006.2668v1, 2010.

\bibitem{Gowers}
W. T. Gowers,
\emph{An infinite Ramsey theorem and some Banach-space dichotomies},
Ann. of Math. (2) 156 (2002), 797--833.

\bibitem{KTJ}
R. A. Komorowski and N. Tomczak-Jaegermann,
\emph{Banach spaces without local unconditional structure},
Israel J. Math. 89 (1995), 205--226.

\end{thebibliography}

\section*{Human Review Recommendation}

Check the finite-class/minimal-class argument and that the use of the
homogeneous Banach-space theorem is exactly at the point where a minimal
representative is isomorphic to all of its own infinite-dimensional subspaces.
The packet should be read as a reduction of Problem 7 to the
\(\ell_2\)-saturated case.

\end{document}
